\chapter{Gibbs Measures and Temperature Decoding}
\label{ch14}

Chapter~13 showed that Boltzmann sampling on combinatorial classes is exactly a Gibbs distribution with energy equal to size. For language models there is a similarly natural energy,
\[
  E(w) = -\log \mu(w),
\]
but now two different temperature operations appear:
\begin{enumerate}
  \item a \emph{global} Gibbs reweighting of complete strings,
  \item the \emph{local} tokenwise temperature scaling used in practice.
\end{enumerate}
The central lesson of the chapter is that these are generally \emph{not the same distribution}.

\section{Sample Space and Energy}

We use the same convention as Chapter~10. Let $\Sigma$ be the content-token vocabulary and let
\[
  \EOS \notin \Sigma
\]
be a stopping symbol. The model samples tokens from $\Sigma \cup \{\EOS\}$, but the actual finite outputs are strings in $\Sigma^*$.

If
\[
  w = w_1\cdots w_n \in \Sigma^*,
\]
then its probability under the autoregressive model is
\[
  \mu(w) = p(\EOS \mid w_1\cdots w_n) \prod_{t=1}^n p(w_t \mid w_1\cdots w_{t-1}).
\]
This is the probability that the model emits exactly the content string $w$ and then stops.

\begin{definition}
The \emph{energy} of a complete output string is
\[
  E(w) = -\log \mu(w).
\]
\end{definition}

The energy is additive along the sampled path because logarithms turn the autoregressive product into a sum of surprisals.

\section{The Global Gibbs Distribution}

For inverse temperature $\beta>0$, define the global Gibbs reweighting by
\[
  \mu_\beta(w) = \frac{\mu(w)^\beta}{Z(\beta)},
  \qquad
  Z(\beta) = \sum_{w \in \Sigma^*} \mu(w)^\beta.
\]
This is a genuine probability distribution whenever the partition function $Z(\beta)$ is finite.

At $\beta=1$, if the original model is tight, then
\[
  Z(1) = \sum_{w \in \Sigma^*} \mu(w) = 1,
\]
so $\mu_1=\mu$.

A useful exact observation is:
\begin{itemize}
  \item for $\beta \ge 1$, one always has $Z(\beta) \le 1$ because $0 \le \mu(w) \le 1$ and raising to a power at least one can only decrease each term;
  \item for $0<\beta<1$, finiteness is no longer automatic and must be checked.
\end{itemize}
So the interesting convergence questions begin on the high-temperature side.

\section{Local Temperature Sampling}

In practical decoding, temperature is applied not to complete-string probabilities but to each next-token distribution separately. For a history $h \in \Sigma^*$, define the local partition factor
\[
  \widetilde Z_\beta(h) = \sum_{a \in \Sigma \cup \{\EOS\}} p(a \mid h)^\beta.
\]
The locally tempered next-token distribution is
\[
  \widetilde p_\beta(a \mid h) = \frac{p(a \mid h)^\beta}{\widetilde Z_\beta(h)}.
\]
Iterating these conditionals produces a new output distribution on $\Sigma^*$, which we denote by $\widetilde\mu_\beta$.

If
\[
  w = w_1\cdots w_n,
\]
then the local-temperature probability of the complete output $w$ is
\[
  \widetilde\mu_\beta(w)
  = \left(\prod_{t=1}^n \frac{p(w_t \mid w_1\cdots w_{t-1})^\beta}{\widetilde Z_\beta(w_1\cdots w_{t-1})}\right)
    \frac{p(\EOS \mid w)^\beta}{\widetilde Z_\beta(w)}.
\]

\section{The Exact Comparison Formula}

Let
\[
  \operatorname{Pref}(w) = \{\varepsilon, w_1, w_1w_2, \dots, w_1\cdots w_n\}
\]
be the set of prefixes encountered while generating $w$, including the full content string just before emitting $\EOS$.

\begin{proposition}
For every complete output string $w \in \Sigma^*$,
\[
  \widetilde\mu_\beta(w) = \frac{\mu(w)^\beta}{\prod_{h \in \operatorname{Pref}(w)} \widetilde Z_\beta(h)}.
\]
Hence, whenever $Z(\beta)<\infty$,
\[
  \frac{\widetilde\mu_\beta(w)}{\mu_\beta(w)}
  = \frac{Z(\beta)}{\prod_{h \in \operatorname{Pref}(w)} \widetilde Z_\beta(h)}.
\]
Therefore
\[
  \widetilde\mu_\beta = \mu_\beta
\]
if and only if the prefix-product
\[
  \prod_{h \in \operatorname{Pref}(w)} \widetilde Z_\beta(h)
\]
is the same for every output string $w$ with positive probability.
\end{proposition}

This is the sharp equality condition. It is stronger than requiring each individual local partition function to be ``prefix-independent'' in isolation, because the full product also depends on how many prefixes a string traverses.

\section{Why the Difference Occurs}

The global Gibbs distribution reweights \emph{complete strings}. The local rule reweights one-step choices and renormalizes immediately after each prefix. That renormalization forgets how much total future mass lies below each branch.

One useful way to say this is to define the future partition factor
\[
  G_\beta(h) = \sum_{u \in \Sigma^*} \mu_h(u)^\beta,
\]
where $\mu_h$ is the continuation distribution from prefix $h$. Then the global Gibbs conditional at prefix $h$ has the form
\[
  \mu_\beta(a \mid h) \propto p(a \mid h)^\beta G_\beta(ha),
\]
with the special case of halting obtained by taking $a=\EOS$ and $G_\beta(h\EOS)=1$.

The local rule keeps the factor $p(a \mid h)^\beta$ but drops the continuation weight $G_\beta(ha)$. That missing subtree mass is exactly where the distortion comes from.

\section{A Fully Specified Example}

Consider the following autoregressive model over content vocabulary $\Sigma=\{a,b,c,d\}$.

At the root:
\[
  p(a \mid \varepsilon)=0.8,
  \qquad
  p(b \mid \varepsilon)=0.2,
  \qquad
  p(\EOS \mid \varepsilon)=0.
\]
After prefix $a$:
\[
  p(\EOS \mid a)=0.9,
  \qquad
  p(c \mid a)=0.1.
\]
After prefix $b$:
\[
  p(\EOS \mid b)=0.5,
  \qquad
  p(d \mid b)=0.5.
\]
After prefixes $ac$ and $bd$, the model halts with probability $1$.

So the four complete outputs are
\[
  a,
  \qquad
  ac,
  \qquad
  b,
  \qquad
  bd,
\]
with probabilities
\[
  \mu(a)=0.72,
  \quad
  \mu(ac)=0.08,
  \quad
  \mu(b)=0.10,
  \quad
  \mu(bd)=0.10.
\]
Take $\beta=1/2$.

\paragraph{Global Gibbs.}
The weights are
\[
  \sqrt{0.72},\; \sqrt{0.08},\; \sqrt{0.10},\; \sqrt{0.10},
\]
whose normalized values are approximately
\[
  \mu_{1/2}(a)\approx 0.481,
  \quad
  \mu_{1/2}(ac)\approx 0.160,
  \quad
  \mu_{1/2}(b)\approx 0.179,
  \quad
  \mu_{1/2}(bd)\approx 0.179.
\]

\paragraph{Local temperature.}
At the root, the tempered distribution on $\{a,b\}$ is proportional to
\[
  \sqrt{0.8},\; \sqrt{0.2},
\]
so
\[
  \widetilde p_{1/2}(a \mid \varepsilon) \approx 0.667,
  \qquad
  \widetilde p_{1/2}(b \mid \varepsilon) \approx 0.333.
\]
At prefix $a$, the local tempered split between halting and continuing is proportional to
\[
  \sqrt{0.9},\; \sqrt{0.1},
\]
so
\[
  \widetilde p_{1/2}(\EOS \mid a) = 0.75,
  \qquad
  \widetilde p_{1/2}(c \mid a)=0.25.
\]
At prefix $b$, the local partition factors are equal, so
\[
  \widetilde p_{1/2}(\EOS \mid b)=\widetilde p_{1/2}(d \mid b)=0.5.
\]
Thus
\[
  \widetilde\mu_{1/2}(a)=0.500,
  \quad
  \widetilde\mu_{1/2}(ac)=0.167,
  \quad
  \widetilde\mu_{1/2}(b)=0.167,
  \quad
  \widetilde\mu_{1/2}(bd)=0.167.
\]
These do not match the global Gibbs probabilities. The low-probability $b$-subtree is over-promoted because the local procedure does not know that its future mass is smaller.

\section{Greedy, High Temperature, and Truncation Heuristics}

Two limits should be distinguished.
Recall that local temperature is usually written as $T$, with $\beta = 1/T$.

\paragraph{Global Gibbs limit.}
As $\beta \to \infty$, the global Gibbs measure concentrates on complete strings that maximize $\mu(w)$ globally.

\paragraph{Local decoding limit.}
As $T \to 0$ in local tokenwise temperature sampling, each next-token distribution concentrates on its local argmax set. This is greedy decoding. Greedy decoding need not coincide with the globally most probable complete string.

At the other extreme, as $\beta \to 0$, the global Gibbs weights flatten. But on a countably infinite set of strings there is no honest uniform distribution to converge to, so ``uniform'' should only be understood heuristically.

Practical heuristics such as top-$k$ and nucleus sampling remain autoregressive procedures: they still define prefix-conditioned next-token distributions. What they break is not the prefix-Markov structure, but the clean interpretation as either the original model or a global Gibbs reweighting of it.

\section*{Preview of Chapter 15}

Chapter~15 studies what happens when the global partition function
\[
  Z(\beta)=\sum_{w \in \Sigma^*} \mu(w)^\beta
\]
changes analytic character as $\beta$ varies. The point of the present chapter is that one must first decide which temperature operation is being studied. Local tokenwise temperature and global Gibbs reweighting are different objects, and any talk of ``phase transitions in temperature decoding'' must say which one it means.

\section*{Further Reading}

\begin{itemize}
  \item \textbf{M. Blondel et al.}, ``Autoregressive language models are secretly energy-based models'' (2025) \cite{BlondelEtAl2025} --- for global energy-based reformulations of autoregressive models.
  \item \textbf{T. Kempton and C. Burrell}, ``Temperature and local normalization in autoregressive models'' (2025) \cite{KemptonBurrell2025} --- for quantitative bounds on the gap between local and global temperature manipulations.
  \item \textbf{P. Flajolet and R. Sedgewick}, \emph{Analytic Combinatorics} (2009) \cite{FlajoletSedgewick2009} --- for the Gibbs/Boltzmann side on the combinatorial end.
\end{itemize}

\section{Exercises}

\begin{exercise}
Verify the ratio formula for $\widetilde\mu_\beta(w)/\mu_\beta(w)$ directly from the definitions.
\end{exercise}

\begin{exercise}
Construct a memoryless model with two content tokens and one EOS token, compute both $\mu_\beta$ and $\widetilde\mu_\beta$, and check that they are generally different when output lengths vary.
\end{exercise}
